\documentclass{article}

% \usepackage[preprint]{neurips_2023} 
\usepackage[margin=1in]{geometry} % Use standard geometry for local compilation

\usepackage[utf8]{inputenc} % allow utf-8 input
\usepackage[T1]{fontenc}    % use 8-bit T1 fonts
\usepackage{hyperref}       % hyperlinks
\usepackage{url}            % simple URL typesetting
\usepackage{booktabs}       % professional-quality tables
\usepackage{amsfonts}       % blackboard math symbols
\usepackage{nicefrac}       % compact symbols for 1/2, etc.
\usepackage{microtype}      % microtypography
\usepackage{xcolor}         % colors
\usepackage{amsmath}
\usepackage{amssymb}
\usepackage{amsthm}
\usepackage{algorithm}
\usepackage{algorithmic}
\usepackage{graphicx}

\newtheorem{theorem}{Theorem}
\newtheorem{lemma}[theorem]{Lemma}
\newtheorem{definition}{Definition}
\newtheorem{assumption}{Assumption}
\newtheorem{proposition}[theorem]{Proposition}

\title{CTAP: Confidence-Guided TEI-Maximizing Adaptive Pruning for Multi-Agent Communication}

\author{
  Anonymous Authors \\
  Affiliation \\
  \texttt{email@domain.edu} \\
}

\begin{document}

\maketitle

\begin{abstract}
Multi-Agent Reinforcement Learning (MARL) relies heavily on communication to foster cooperation. However, unconstrained communication often leads to information redundancy and high bandwidth costs. Existing pruning methods either rely on static thresholds or instantaneous attention weights, failing to evaluate the long-term, true contribution of communication channels. In this paper, we propose Confidence-guided TEI-maximizing Adaptive Pruning (CTAP), a novel online communication topology optimization framework. CTAP online-estimates the edge-level Shapley value of each communication link to measure its exact contribution. To prevent erroneous pruning caused by the stochasticity of MARL environments, we introduce a confidence-bound mechanism equivalent to sequential hypothesis testing, and theoretically prove its effectiveness in bounding mispruning probabilities. Furthermore, we optimize a Topology Efficiency Index (TEI) representing a Lagrangian relaxation of the communication budget, utilizing a zeroth-order stochastic approximation method to dynamically adjust the pruning threshold. We theoretically prove the optimal elasticity condition for TEI and the topological stability under submodularity. Extensive experiments on SMACv2 and MPE benchmarks demonstrate that CTAP significantly reduces communication overhead while maintaining or even improving the cooperative performance compared to state-of-the-art baselines.
\end{abstract}

\section{Introduction}
In Multi-Agent Reinforcement Learning (MARL), communication allows agents to share local observations and coordinate actions. While fully connected communication graphs ensure maximum information exchange, they impose severe bandwidth overheads and are prone to the ``lazy agent'' problem, where agents become overly reliant on others' messages. 

To address this, recent studies focus on learning sparse communication topologies. However, determining \textit{which} message to drop remains a challenge. Attention-based methods (e.g., TarMAC) use attention weights as a proxy for importance, but attention does not strictly align with causal contribution to the global reward. Game-theoretic measures, specifically the Shapley value, provide a rigorous way to distribute the total reward among communicating edges. Yet, exact Shapley computation is NP-hard, and its empirical estimation in highly stochastic environments suffers from high variance, leading to potential mispruning of beneficial edges.

We propose CTAP, an online framework that leverages the edge-level Shapley value with a theoretically grounded confidence evaluation to prune detrimental or redundant communication edges. Additionally, CTAP introduces a dynamically adaptive threshold mechanism that maximizes the Topology Efficiency Index (TEI)---the ratio of task success rate to communication density.

\section{Methodology and Theoretical Analysis}
Our framework evaluates the contribution of each directed communication edge $e=(i,j)$ using its online-estimated Shapley value $\phi_e$. To handle environmental stochasticity and dynamic topology adjustments, we introduce the Confidence-Guided Pruning and TEI-Maximizing Adaptive Threshold, both supported by rigorous theoretical guarantees.
\begin{figure*}[t]
    \centering
    \includegraphics[width=1\textwidth]{figures/CTAP_Flowchart.pdf}
    \caption{Overview of the CTAP framework. Starting from a fully connected communication graph \textbf{(1)}, edge-level Shapley values are estimated online every $K$ epochs \textbf{(2)}, assessed via temporal confidence modeling and TEI-driven threshold optimization \textbf{(3)}, and used to update a binary communication mask that produces the pruned topology \textbf{(4)}. Green solid edges are retained; red dashed edges are pruned.}
    \label{fig:CTAP_Flowchart}
\end{figure*}

As illustrated in Figure~\ref{fig:CTAP_Flowchart}, CTAP operates as a four-stage closed-loop pipeline interleaving regular policy training with periodic topology refinement.
\textbf{Stage 1 (Multi-Agent Environment).} Agents are initialized with a fully connected directed communication graph, where each edge $(i \to j)$ represents a dedicated message-passing channel from agent $i$ to agent $j$.
\textbf{Stage 2 (Online Shapley Value Estimation).} Every $K$ training epochs, the current policy is frozen and Monte Carlo rollouts are executed to estimate the edge-level Shapley value $\varphi(i \to j) = \mathbb{E}[V(S \cup \{i \to j\})] - V(S)$ for each directed edge. The resulting Shapley matrix explicitly quantifies the marginal causal contribution of each communication channel to the team reward.
\textbf{Stage 3 (Confidence \& TEI Assessment).} The estimated Shapley values are evaluated through two complementary mechanisms operating in parallel. \textit{Temporal Confidence Modeling} accumulates Shapley estimates over a sliding window of size $W$ and computes the confidence score $\text{Conf}(i,j) = |\mu(\varphi)| / (1 + \sigma(\varphi))$, which identifies edges that are consistently detrimental with high statistical confidence. \textit{TEI-Driven Threshold Optimization} monitors the Task Efficiency Index $\text{TEI} = \text{Success Rate} / \text{Communication Density}$ and adapts the pruning threshold $\theta(t)$ via zeroth-order gradient ascent: the threshold is tightened when TEI improves and relaxed when TEI degrades, continuously discovering the Pareto-optimal communication budget.
\textbf{Stage 4 (Adaptive Graph Pruning).} A binary communication mask is updated by pruning edge $(i \to j)$ if and only if it simultaneously satisfies: $\text{Conf}(i,j) \ge 0$ (stable negative contributor) \textit{and} $\mu(\varphi_{ij}) < \theta(t)$ (below the adaptive threshold). The dual-condition gate prevents erroneous pruning caused by transient noise, ensuring only persistently harmful channels are removed. The resulting sparse topology is applied to subsequent training, and the cycle repeats until convergence to a CTAP-stable graph. The complete procedure is summarized in Algorithm~\ref{algo:CTAP}. In the following subsections, we provide the formal theoretical foundation for each stage of this pipeline.

\begin{algorithm}[H]
\caption{CTAP: Confidence-Guided TEI-Maximizing Adaptive Pruning}
\label{algo:CTAP}
\begin{algorithmic}[1]
\REQUIRE Policy parameters $\boldsymbol{\theta}$; fully connected graph $G_0=(A,E)$ with $n$ agents; mask $\mathbf{M}\in\{0,1\}^{n\times n}$ (all ones); initial threshold $\theta(0)$; window size $W$; confidence threshold $p_c$; evaluation interval $K$; sequences $\{c_t\}$, $\{\alpha_t\}$.
\ENSURE Converged policy $\boldsymbol{\theta}^*$ and CTAP-stable communication graph $G^*$.
\FOR{epoch $t = 1, 2, \ldots$ \textbf{until convergence}}
    \IF{$t \bmod K = 0$}
        \STATE $\triangleright$ \textit{Stage 1--2: Policy Rollout and Online Shapley Estimation}
        \STATE Freeze policy parameters $\boldsymbol{\theta}$.
        \FORALL{directed edge $e = (i \to j) \in E$}
            \STATE Sample $m$ MC episodes with random subsets $S \subseteq E \setminus \{e\}$.
            \STATE $\Delta_k(e) \leftarrow V(S_k \cup \{e\}) - V(S_k)$,\; $k = 1, \ldots, m$.
            \STATE $\hat{\varphi}_e \leftarrow \frac{1}{m}\sum_{k=1}^{m} \Delta_k(e)$;\; append to $\mathcal{W}_e$, drop oldest if $|\mathcal{W}_e| > W$.
        \ENDFOR
        \STATE Unfreeze policy parameters $\boldsymbol{\theta}$.
    \ENDIF
    \STATE $\triangleright$ \textit{Stage 3a: Temporal Confidence Modeling}
    \FORALL{edge $e = (i \to j) \in E$}
        \STATE $\mu_e \leftarrow \mathrm{mean}(\mathcal{W}_e)$,\; $\sigma_e \leftarrow \mathrm{std}(\mathcal{W}_e)$,\; $\mathrm{Conf}(e) \leftarrow |\mu_e| / (1 + \sigma_e)$.
        \STATE $\hat{p}_e \leftarrow |\{v \in \mathcal{W}_e : v < 0\}| / W$.
    \ENDFOR
    \STATE $\triangleright$ \textit{Stage 3b: TEI-Driven Threshold Optimization}
    \STATE Evaluate $\widehat{\mathrm{TEI}}(\theta(t) + c_t)$ and $\widehat{\mathrm{TEI}}(\theta(t) - c_t)$ from rollout statistics.
    \STATE $\hat{g}_t \leftarrow \bigl[\widehat{\mathrm{TEI}}(\theta(t)+c_t) - \widehat{\mathrm{TEI}}(\theta(t)-c_t)\bigr] / (2c_t)$.
    \STATE $\theta(t+1) \leftarrow \Pi_{[\theta_{\min},\theta_{\max}]}\bigl[\theta(t) + \alpha_t \hat{g}_t\bigr]$.
    \STATE $\triangleright$ \textit{Stage 4: Adaptive Graph Pruning}
    \FORALL{edge $e = (i \to j) \in E$}
        \IF{$\hat{p}_e \ge p_c$ \textbf{and} $\mu_e < \theta(t+1)$}
            \STATE $\mathbf{M}_{ij} \leftarrow 0$ \COMMENT{Prune: stable negative contributor below threshold}
        \ELSE
            \STATE $\mathbf{M}_{ij} \leftarrow 1$ \COMMENT{Retain}
        \ENDIF
    \ENDFOR
    \STATE $G_t \leftarrow \{e \in E : \mathbf{M}_{ij} = 1\}$;\; update $\boldsymbol{\theta}$ via policy gradient under $G_t$.
\ENDFOR
\RETURN $\boldsymbol{\theta}^*$,\; $G^*$.
\end{algorithmic}
\end{algorithm}

\subsection{Axiomatic Foundation and Sample Complexity of Edge-Level Shapley Value}
Let $E$ be the set of all possible directed communication edges among $n$ agents. The edge-level Shapley value $\phi(i \to j)$ evaluates the expected marginal contribution of the edge to the coalition:
\begin{equation}
    \phi(i \to j) = \sum_{S \subseteq E \setminus \{(i,j)\}} \frac{|S|!(|E|-|S|-1)!}{|E|!} \left[ V(S \cup \{(i,j)\}) - V(S) \right]
\end{equation}
where $V(S)$ is the expected team reward when only the edges in subset $S$ are activated. The Shapley value is uniquely characterized by four axioms: Efficiency, Symmetry, Null Player, and Linearity. Therefore, the pruning criterion in CTAP is game-theoretically \textit{fair}. 

Since exact computation is intractable, CTAP employs Monte Carlo sampling. By keeping the policy fixed during the evaluation phase, the sampled episodes are treated as independent and identically distributed (i.i.d.). We establish the sample complexity of this estimation.

\begin{theorem}[Sample Complexity of Online Estimation]
Assuming the team reward is bounded $V(\cdot) \in [R_{\min}, R_{\max}]$ and evaluations are i.i.d., estimating $\hat{\phi}(i \to j)$ via $m$ Monte Carlo episodes yields an estimation error bounded with high probability $1-\delta$:
\begin{equation}
    \mathbb{P}\left(\max_{e \in E}|\hat{\phi}_e - \phi_e| \le (R_{\max} - R_{\min}) \sqrt{\frac{2\ln(2n^2/\delta)}{m}}\right) \ge 1 - \delta
\end{equation}
\end{theorem}
\begin{proof}
Let $X_k$ denote the marginal contribution of edge $e$ evaluated in the $k$-th independent sampled episode for $k \in \{1, \dots, m\}$. The estimation is given by the empirical mean $\hat{\phi}_e = \frac{1}{m}\sum_{k=1}^m X_k$.

Given the reward bounds $V(\cdot) \in [R_{\min}, R_{\max}]$, the marginal contribution is bounded as:
\begin{equation}
    X_k = V(S \cup \{e\}) - V(S) \in [-(R_{\max} - R_{\min}), R_{\max} - R_{\min}]
\end{equation}
yielding a maximum interval length $C = 2(R_{\max} - R_{\min})$.

Applying Hoeffding's Inequality for the sum of bounded independent random variables, for any $\varepsilon > 0$:
\begin{equation}
    \mathbb{P}\left(|\hat{\phi}_e - \phi_e| \ge \varepsilon\right) \overset{\text{Hoeffding}}{\le} 2 \exp\left(-\frac{2m^2\varepsilon^2}{\sum_{k=1}^m C^2}\right) \overset{\text{Substitute } C}{=} 2 \exp\left(-\frac{m\varepsilon^2}{2(R_{\max} - R_{\min})^2}\right)
\end{equation}

To ensure the error bounds hold globally across all potential edges simultaneously, we apply the Union Bound over at most $n(n-1) \le n^2$ directed edges:
\begin{align}
    \mathbb{P}\left(\exists e \in E : |\hat{\phi}_e - \phi_e| \ge \varepsilon\right) 
    &\overset{\text{Union Bound}}{\le} \sum_{e \in E} \mathbb{P}\left(|\hat{\phi}_e - \phi_e| \ge \varepsilon\right) \\
    &\overset{|E| \le n^2}{\le} 2n^2 \exp\left(-\frac{m\varepsilon^2}{2(R_{\max} - R_{\min})^2}\right)
\end{align}

We equate this upper bound to the failure probability $\delta$:
\begin{equation}
    2n^2 \exp\left(-\frac{m\varepsilon^2}{2(R_{\max} - R_{\min})^2}\right) \overset{\text{set}}{=} \delta
\end{equation}

Solving for the error margin $\varepsilon$ yields step-by-step:
\begin{align}
    \exp\left(-\frac{m\varepsilon^2}{2(R_{\max} - R_{\min})^2}\right) &\overset{\text{Divide}}{=} \frac{\delta}{2n^2} \\
    \frac{m\varepsilon^2}{2(R_{\max} - R_{\min})^2} &\overset{\text{Logarithm}}{=} \ln\left(\frac{2n^2}{\delta}\right) \\
    \varepsilon &\overset{\text{Solve for } \varepsilon}{=} (R_{\max} - R_{\min}) \sqrt{\frac{2\ln(2n^2/\delta)}{m}}
\end{align}
This explicit derivation demonstrates that achieving a maximum global error $\varepsilon$ requires a sample complexity $m \in \mathcal{O}\left(\frac{(R_{\max}-R_{\min})^2\ln(n^2/\delta)}{\varepsilon^2}\right)$.
\end{proof}

\subsection{Confidence-Guided Pruning via Sequential Hypothesis Testing}
In stochastic environments, an instantaneous estimation $\hat{\phi}_{e, t} < 0$ might just be noise. CTAP maintains a sliding window of size $W$ to compute a confidence score $\mathcal{C}_e = \frac{|\bar{\phi}_e|}{1 + \hat{\sigma}_e}$. This mechanism can be formalized as a sequential hypothesis test.

Consider testing $H_0: \phi(i \to j) \ge 0$ (Beneficial Edge) versus $H_1: \phi(i \to j) < 0$ (Detrimental Edge). Under $H_0$, assuming an unbiased estimator, the probability of observing a negative estimation in a single episode is at most $1/2$. Over a window $W$, the number of negative observations $k$ follows a Binomial distribution $\text{Bin}(W, 0.5)$. CTAP requires the proportion of negative observations to exceed a confidence threshold $p_c$. The significance level $\alpha$ (probability of Type I error, i.e., mispruning) is:
\begin{equation}
    \alpha = \mathbb{P}\left(\text{Bin}(W, 0.5) \ge \lceil p_c \cdot W \rceil\right)
\end{equation}
This explicitly provides the theoretical foundation for selecting hyperparameters $W$ and $p_c$ to bound false discovery rates.

\begin{theorem}[Exponential Bound on Mispruning Probability]
Let the true expected Shapley value of edge $e$ be $\mu_e^* = \mathbb{E}[\phi_e]$, and its true variance be $\sigma_e^2$. Since the reward is bounded, the estimations $\hat{\phi}_e^{(k)}$ strictly satisfy the sub-Gaussian property with a variance proxy parameter proportional to $\sigma_e^2$. Assuming sufficient mixing such that evaluations across the sliding window $W$ are approximately independent, we prune the edge when the sample mean $\bar{\phi}_e$ is negative. Using the proposed confidence score $\mathcal{C}_e = \frac{|\bar{\phi}_e|}{1 + \hat{\sigma}_e}$, the empirical probability of a false negative estimation (mispruning) decays \textit{exponentially}:
\begin{equation}
    \mathbb{P}(\text{Mispruning}) \le \exp\left(- \frac{W}{2} \mathcal{C}_e^2 \right)
\end{equation}
\end{theorem}

\begin{proof}
Let $\bar{\phi}_e = \frac{1}{W} \sum_{k=1}^W \hat{\phi}_{e}^{(k)}$. Since the bounded marginal contributions guarantee the estimations are sub-Gaussian, we apply the Hoeffding-type concentration inequality for the sample mean. For any $\varepsilon > 0$:
\begin{equation}
    \mathbb{P}(|\bar{\phi}_e - \mu_e^*| \ge \varepsilon) \overset{\text{Sub-Gaussian}}{\le} \exp\left(-\frac{W \varepsilon^2}{2\sigma_e^2}\right)
\end{equation}
Assume the edge is fundamentally beneficial, i.e., $\mu_e^* \ge 0$. The event $\bar{\phi}_e < 0$ implies that the sample mean has deviated from the true mean by at least $\mu_e^*$. Substituting $\varepsilon = \mu_e^*$, we derive step-by-step:
\begin{align}
    \mathbb{P}(\bar{\phi}_e < 0 \mid \mu_e^* \ge 0) &\overset{\text{Deviation}}{\le} \mathbb{P}(|\bar{\phi}_e - \mu_e^*| \ge \mu_e^*) \\
    &\overset{\text{Substitute } \varepsilon}{\le} \exp\left(-\frac{W (\mu_e^*)^2}{2\sigma_e^2}\right)
\end{align}
In practice, the unknown true signal-to-noise ratio $\frac{\mu_e^*}{\sigma_e}$ is empirically lower-bounded by our confidence score $\mathcal{C}_e$:
\begin{equation}
    \frac{\mu_e^*}{\sigma_e} \approx \frac{|\bar{\phi}_e|}{\hat{\sigma}_e} > \frac{|\bar{\phi}_e|}{1 + \hat{\sigma}_e} = \mathcal{C}_e
\end{equation}
Substituting this empirical proxy into the bound yields:
\begin{equation}
    \mathbb{P}(\text{Mispruning}) \overset{\text{Apply } \mathcal{C}_e}{\lesssim} \exp\left(- \frac{W}{2} \mathcal{C}_e^2 \right)
\end{equation}
This tightening demonstrates that as the confidence score $\mathcal{C}_e$ increases, the probability of mistakenly pruning a beneficial communication edge decays at a rigorous exponential rate, validating the safety of our confidence-guided mechanism.
\end{proof}

\subsection{TEI-Maximizing Adaptive Threshold and Lagrangian Relaxation}
Instead of relying on a static pruning budget $k$, which requires solving the constrained problem $\max_G R(G) \text{ s.t. } |E(G)| \le k$, CTAP optimizes the Task Efficiency Index (TEI), defined as $\text{TEI}(\tau) = \frac{R(\tau)}{\rho(\tau) + \varepsilon}$, where $R(\tau)$ is the success rate and $\rho(\tau)$ is the communication density. 

Maximizing TEI conceptually mirrors finding the tangent point on the Pareto frontier of the $R-\rho$ plane, dynamically discovering the optimal Lagrangian multiplier $\lambda$ in the relaxed objective $\max_G R(G) - \lambda |E(G)|$.

\begin{assumption}[Continuous Relaxation of Topology]
While the pruning threshold $\tau$ strictly induces a discrete topology graph $G(\tau)$, making $R(\tau)$ and $\rho(\tau)$ piecewise constant, we assume for the purpose of gradient-based optimization that $R(\tau)$ and $\rho(\tau)$ can be approximated by continuously differentiable surrogate functions.
\end{assumption}

\begin{proposition}[TEI Optimality Condition]
Under the continuous relaxation assumption, at the optimal threshold $\tau^*$ (assuming $\varepsilon$ is negligible), the necessary condition for TEI optimality reveals that the elasticity of the success rate equates the elasticity of communication density.
\end{proposition}

\begin{proof}
Let the objective be maximizing $f(\tau) = \frac{R(\tau)}{\rho(\tau)}$. Taking the derivative with respect to $\tau$ and setting it to zero yields:
\begin{equation}
    \frac{d}{d\tau} \left( \frac{R(\tau)}{\rho(\tau)} \right) \overset{\text{Quotient Rule}}{=} \frac{R'(\tau)\rho(\tau) - R(\tau)\rho'(\tau)}{\rho(\tau)^2} \overset{\text{set}}{=} 0
\end{equation}
Since $\rho(\tau) > 0$, the numerator must be zero, which rearranging gives step-by-step:
\begin{align}
    R'(\tau^*)\rho(\tau^*) - R(\tau^*)\rho'(\tau^*) &\overset{\text{Numerator}}{=} 0 \\
    R'(\tau^*)\rho(\tau^*) &\overset{\text{Rearrange}}{=} R(\tau^*)\rho'(\tau^*) \\
    \frac{R'(\tau^*)}{R(\tau^*)} &\overset{\text{Divide}}{=} \frac{\rho'(\tau^*)}{\rho(\tau^*)}
\end{align}
This indicates that at the global optimal threshold, the marginal relative gain in success rate exactly matches the marginal relative reduction in communication density.
\end{proof}

To reach $\tau^*$, CTAP employs a zeroth-order stochastic approximation method based on the Kiefer-Wolfowitz (KW) algorithm. Since the exact gradient is unavailable, we estimate it using finite differences with a perturbation sequence $c_t$.

\begin{theorem}[Convergence of Adaptive Threshold under Biased Gradient]
Let the objective function $J(\tau) = \mathbb{E}[\text{TEI}(\tau)]$ be $L$-smooth. The threshold update rule is $\tau_{t+1} = \Pi_{[\tau_{\min}, \tau_{\max}]} \left[ \tau_t + \alpha_t \hat{\nabla} J(\tau_t) \right]$, where the finite difference gradient estimator $\hat{\nabla} J(\tau_t) = \frac{\hat{J}(\tau_t + c_t) - \hat{J}(\tau_t - c_t)}{2c_t}$.
Due to the non-linearity of $J$, this estimator is biased: $\mathbb{E}[\hat{\nabla} J(\tau_t)] = \nabla J(\tau_t) + \mathcal{O}(c_t^2)$.
If the step size $\alpha_t$ and perturbation $c_t$ satisfy the standard KW conditions: $\sum \alpha_t = \infty$, $\sum \alpha_t c_t < \infty$, and $\sum (\alpha_t / c_t)^2 < \infty$, with $c_t \to 0$, then the sequence $\{\tau_t\}$ converges almost surely to a stationary point of $J(\tau)$.
\end{theorem}

\begin{proof}
Let $b(\tau_t) = \mathbb{E}[\hat{\nabla} J(\tau_t) \mid \mathcal{F}_t] - \nabla J(\tau_t)$ be the estimation bias. By Taylor expansion of $J(\tau \pm c_t)$, the bias is bounded by $\mathcal{O}(c_t^2)$ due to $L$-smoothness. Let the variance of the estimator be bounded by $\mathcal{O}(1/c_t^2)$.

By taking the Taylor expansion of $J(\tau_{t+1})$ around $\tau_t$:
\begin{equation}
    J(\tau_{t+1}) \overset{\text{Taylor}}{\ge} J(\tau_t) + \langle \nabla J(\tau_t), \tau_{t+1} - \tau_t \rangle - \frac{L}{2} \|\tau_{t+1} - \tau_t\|^2
\end{equation}
Taking the expectation conditioned on the filtration $\mathcal{F}_t$ up to epoch $t$:
\begin{align}
    \mathbb{E}[J(\tau_{t+1}) \mid \mathcal{F}_t] &\overset{\text{Update Rule}}{\ge} J(\tau_t) + \alpha_t \langle \nabla J(\tau_t), \nabla J(\tau_t) + b(\tau_t) \rangle - \frac{L \alpha_t^2}{2} \mathbb{E}[\|\hat{\nabla} J(\tau_t)\|^2 \mid \mathcal{F}_t] \\
    &\overset{\text{Bounded Var}}{\ge} J(\tau_t) + \alpha_t \|\nabla J(\tau_t)\|^2 + \alpha_t \langle \nabla J(\tau_t), b(\tau_t) \rangle - \mathcal{O}\left( \frac{\alpha_t^2}{c_t^2} \right)
\end{align}
Using Cauchy-Schwarz on the bias term $\alpha_t \langle \nabla J(\tau_t), b(\tau_t) \rangle \ge - \alpha_t \|\nabla J(\tau_t)\| \|b(\tau_t)\|$, and noting that $\|b(\tau_t)\| \le M c_t^2 \le M c_t$ for small $c_t$, the negative impacts are bounded.
Rearranging the terms and summing over $t=1$ to $T$:
\begin{align}
    \sum_{t=1}^T \alpha_t \|\nabla J(\tau_t)\|^2 &\overset{\text{Sum \& Telescoping}}{\le} \mathbb{E}[J(\tau_{T+1})] - J(\tau_1) + \sum_{t=1}^T \alpha_t \|\nabla J(\tau_t)\| \|b(\tau_t)\| + \mathcal{O}\left( \sum_{t=1}^T \frac{\alpha_t^2}{c_t^2} \right)
\end{align}
Since $J(\tau)$ is bounded on $[\tau_{\min}, \tau_{\max}]$, and by the KW conditions $\sum \alpha_t c_t < \infty$ (bounding the bias sum) and $\sum (\alpha_t / c_t)^2 < \infty$ (bounding the variance sum), the right-hand side is bounded by a constant $M' < \infty$. Given $\sum \alpha_t = \infty$, the series on the left can only converge if:
\begin{equation}
    \liminf_{t \to \infty} \|\nabla J(\tau_t)\| \overset{\text{Limit}}{=} 0 \quad \text{almost surely.}
\end{equation}
Thus, despite the biased finite-difference gradient, the adaptive threshold strictly converges to a local optimum of the TEI objective.
\end{proof}

\subsection{Fixed-Point Analysis of Topology Stability}
Since the Shapley value $\phi_G(e)$ depends on the current topology $G$, the adaptive pruning forms a feedback loop. A topology $G^*$ is defined as \textbf{CTAP-stable} at threshold $\tau^*$ if and only if:
\begin{equation}
    \forall e \in G^*, \phi_{G^*}(e) \ge \tau^* \quad \text{and} \quad \forall e \notin G^*, \phi_{G^*}(e) < \tau^*
\end{equation}

\begin{assumption}[Submodularity of Information Aggregation]
We assume the value function $V(S)$ is monotonically submodular, representing the principle of diminishing marginal returns for information in cooperative MARL. While specific synergistic interactions might exhibit supermodularity, the general aggregation of shared observations typically provides decreasing incremental benefits as the network approaches full connectivity.
\end{assumption}

\begin{theorem}[Monotonic Convergence under Submodularity]
Under the assumption of submodularity, the edge-level Shapley value $\phi_G(e)$ is monotonically non-increasing with respect to the graph $G$. Consequently, the CTAP iterative pruning process is monotonically contractive and guaranteed to converge to a stable fixed point without oscillation.
\end{theorem}

\begin{proof}
Submodularity of the value function $V(S)$ implies that for any subset of edges $A \subseteq B \subset E$ and any edge $e \notin B$, the marginal contribution satisfies:
\begin{equation}
    V(A \cup \{e\}) - V(A) \overset{\text{Submodularity}}{\ge} V(B \cup \{e\}) - V(B)
\end{equation}
The Shapley value $\phi_G(e)$ is a convex combination of these marginal contributions evaluated over subsets $S \subseteq G \setminus \{e\}$. Consequently, since every marginal contribution decreases as the underlying graph grows, the expectation over these contributions must also follow the same inequality. Therefore, if $A \subseteq B$, we have step-by-step:
\begin{align}
    \mathbb{E}_{S \subseteq A}[V(S \cup \{e\}) - V(S)] &\overset{\text{Marginal}}{\ge} \mathbb{E}_{S \subseteq B}[V(S \cup \{e\}) - V(S)] \\
    \phi_A(e) &\overset{\text{Shapley}}{\ge} \phi_B(e)
\end{align}

In the CTAP pruning process, edges whose Shapley values are less than the threshold $\tau^*$ are removed. Let the initial graph be $G_0 = E$, and the graph at iteration $t$ be $G_t$. Since $G_1 \subset G_0$, by submodularity, the Shapley values of the surviving edges in $G_1$ must satisfy:
\begin{equation}
    \phi_{G_1}(e) \overset{G_1 \subset G_0}{\ge} \phi_{G_0}(e)
\end{equation}
Since the surviving edges already satisfied $\phi_{G_0}(e) \ge \tau^*$, they will strictly satisfy:
\begin{equation}
    \phi_{G_1}(e) \ge \phi_{G_0}(e) \overset{\text{Surviving Edge}}{\ge} \tau^*
\end{equation}
Thus, no previously preserved edge will ever fall below the threshold and be pruned in a subsequent step solely due to the pruning of other edges.

Similarly, an edge $e$ pruned because $\phi_{G_0}(e) < \tau^*$ cannot be evaluated again unless we add it back, but our pruning is subtractive. Therefore, the sequence of graphs forms a monotonic chain:
\begin{equation}
    G_0 \overset{\text{Subtractive Pruning}}{\supseteq} G_1 \supseteq G_2 \dots
\end{equation}
Since the graph is finite, the sequence must contract to a stable fixed point $G^*$ within a finite number of iterations, completely eliminating the possibility of cyclic oscillations.
\end{proof}

\section{Experiments}
\subsection{Experimental Setup}
\textbf{Environments:}
We evaluate CTAP on two popular MARL benchmarks:
\begin{itemize}
    \item \textbf{StarCraft Multi-Agent Challenge v2 (SMACv2):} A complex, stochastic micro-management benchmark. We utilize dynamically generated maps (e.g., \texttt{10gen\_terran} restricted to 5v5 scenarios).
    \item \textbf{Multi-Agent Particle Environments (MPE):} Continuous observation tasks. We focus on \texttt{simple\_spread} and \texttt{simple\_speaker\_listener}.
\end{itemize}

\textbf{Baselines:}
We compare CTAP against state-of-the-art communication architectures with their native pruning or gating mechanisms:
\begin{enumerate}
    \item \textbf{MAGIC:} Graph-attention based communication.
    \item \textbf{CommNet:} Continuous mean-pooling broadcast.
    \item \textbf{IC3Net:} Hard attention with individual gating.
    \item \textbf{GAComm:} Graph-attention with edge penalty.
    \item \textbf{TarMAC:} Targeted multi-agent communication via soft attention.
\end{enumerate}

\subsection{Performance and Communication Trade-offs (Placeholder)}
\textit{[In this section, insert tables/figures comparing the Success Rate and Communication Density across baselines vs. CTAP. Analyze how CTAP achieves comparable or higher success rates while dropping communication bandwidth significantly.]}

\subsection{Ablation Study (Placeholder)}
\textit{[In this section, compare the `full' CTAP mode against `conf\_only' and `tei\_only' to demonstrate the necessity of both the confidence modeling and the dynamic TEI optimization.]}

\section{Conclusion}
In this paper, we introduced CTAP, a rigorous and theoretically grounded framework for pruning communication topologies in MARL. By leveraging confidence-bounded Shapley values, TEI-maximizing zeroth-order optimization, and establishing the optimal elasticity conditions and submodular fixed-point stability, CTAP addresses the non-stationarity and overhead challenges in multi-agent communication. Empirical results validate its superiority in maintaining high task performance with minimal communication bandwidth.

\bibliographystyle{plain}
\bibliography{references} % Placeholder for references

\end{document}